% Subsection: which functional form describes the compute -> ROC AUC frontier, and is its ceiling identified?
% All numbers come from EffingScaling/playground/sigmoid_fit.py (round 5, after the 2026-10-06 fixes) run on
% data/nanotabpfn_data_preprocessed.csv as regenerated on 2026-10-06 (59,360 rows):
%   python sigmoid_fit.py --round 5 --frontier all --bootstrap 100           -> figures/.../functional_forms/round_5_all/
%   python sigmoid_fit.py --round 5 --frontier all --topk 5 --bootstrap 100  -> round_5_all_top5/
%   python sigmoid_fit.py --round 5 --frontier all --seed-weights            -> round_5_all_seedw/
%   python sigmoid_fit.py --round 5 --frontier big_prior --no-loro           -> round_5/
% Horizon-split ladder errors and the ceiling-drift table were computed from the same fits (see appendix).
% Details are in functional_forms_appendix.tex. Requires amsmath and booktabs.
% Assumed bibtex keys (arXiv ids): hoffmann2022chinchilla (2203.15556), alabdulmohsin2022revisiting (2209.06640),
% caballero2023bnsl (2210.14891), rosenfeld2020constructive (1909.12673), arora2025bayesian (2410.16531),
% viering2022shape (2103.10948), owen2024predictable (2401.04757), ruan2024observational (2405.10938),
% gadre2024overtraining (2403.08540), brandfonbrener2024losstoloss (2411.12925).

\subsection{Which functional form describes the ROC AUC frontier?}
\label{sec:functional-forms}

Real-data ROC AUC is bounded, and our compute frontier for it is far from linear in $\log C$: it rises from 0.61 at
$1.8\times10^{11}$ FLOPs to 0.87 at $3\times10^{16}$ FLOPs and visibly flattens. Scaling laws for such metrics usually add
an irreducible term, a sigmoid, or a link from a loss law. We compare 31 such forms on how well they extrapolate, and ask
whether the data determine the ceiling, the ROC AUC that the forms predict for unlimited compute.

\paragraph{Frontier.}
At each of 22 compute levels, from $10^{11.25}$ to $10^{16.5}$ FLOPs in quarter decades, we take the best of up to
6{,}000 runs, and keep a level only if it improves on all smaller ones. The 22 points come from 18 runs, with at most four
checkpoints per run. An earlier version of this analysis used only runs with the largest prior (16 features and 400
rows). That restriction removes the whole $10^{13}$-FLOP stage, whose largest prior has 8 features. Below $10^{13.5}$
FLOPs the restricted frontier therefore trails the full one by up to 0.20 ROC AUC, and above $10^{14}$ FLOPs the two
agree to within 0.006. Its low-compute plateau is a consequence of which runs it admits, not of how performance scales.
On it, no form extrapolates (Appendix~\ref{app:functional-forms}).

\paragraph{Forms.}
Apart from two references without one, every form has a ceiling. We fit power laws with an irreducible term on the gap $1-\text{AUC}$
\citep{hoffmann2022chinchilla} and on $\operatorname{logit}(\text{AUC})$, a logistic in $\log C$, M4
\citep{alabdulmohsin2022revisiting}, a bounded power law \citep{arora2025bayesian,viering2022shape}, the envelope of
\citet{rosenfeld2020constructive}, which bounds the gap by its chance level, and a Weibull approach to the ceiling. We
also fit BNSL with one and two breaks \citep{caballero2023bnsl}, four reduced two-break variants, and power laws with
transient steps. Two-stage forms first fit a power law with an irreducible term to a loss, and then map that loss to ROC
AUC with a logistic \citep{owen2024predictable,ruan2024observational} or exponential \citep{gadre2024overtraining} link.
The loss is the synthetic validation loss, the real-data NLL, or the real-data NLL predicted from the synthetic loss by
a loss-to-loss law \citep{brandfonbrener2024losstoloss}. Table~\ref{tab:forms-appendix} lists all forms.

\paragraph{Evaluation.}
We rank forms by extrapolation. For every level that leaves at least three points above it, we fit on the points up to
that level and measure the mean absolute error (MAE) on the points above. We report this error separately for long
(fit up to $10^{13}$ FLOPs, on 4--8 points), medium and short horizons. A leave-one-run-out test checks interpolation
without sharing a learning curve between training and test points. For the ceiling we report a profile, the range of
ceilings whose fit loss stays within 10\% of the optimum, and a 95\% interval from bootstrapping the 6{,}075 runs and
rebuilding the frontier from each resample. Two details of this protocol matter more than any choice of form. Bounds that
force the ceiling above the best observed ROC AUC must use only the training points. Holdout refits must not start from
the parameters fitted on all points. Each mistake leaks test points into the refit, and together they made a nine-parameter
two-break BNSL look like the best extrapolator in an earlier version of this analysis.

\begin{table}[t]
  \centering
  \caption{Extrapolation error and ceiling of selected forms on the compute--ROC AUC frontier (22 points). MAE in ROC
  AUC, averaged over the cuts of each horizon: long, fit up to $10^{13}$ FLOPs (5 cuts); medium, up to $10^{13.25}$ to
  $10^{14.5}$ (6 cuts); short, up to $10^{14.75}$ to $10^{15.75}$ (5 cuts). LORO: leave one run out. Ceiling: fit on all
  points, with its 95\% run-bootstrap interval; $^{\dagger}$the profile does not bound the ceiling. $k$: free
  parameters. The full table is in the appendix.}
  \label{tab:forms}
  \footnotesize
  \setlength{\tabcolsep}{4pt}
  \begin{tabular}{lccccccc}
    \toprule
    & & \multicolumn{3}{c}{Extrapolation MAE} & & \multicolumn{2}{c}{Ceiling} \\
    \cmidrule(lr){3-5}\cmidrule(lr){7-8}
    Form & $k$ & Long & Medium & Short & LORO & Fit & Bootstrap \\
    \midrule
    Two-stage via real NLL, logistic link   & 4 & 0.022 & 0.014 & 0.013 & 0.006 & 0.897 & [0.883, 0.901] \\
    Power law, logit space                  & 3 & 0.024 & 0.009 & 0.009 & 0.005 & 0.896 & [0.885, 0.899] \\
    Two-stage via synthetic loss, logistic & 4 & 0.026 & 0.012 & 0.013 & 0.006 & 0.909 & [0.883, 0.921] \\
    M4                                      & 4 & 0.032 & 0.015 & 0.011 & 0.007 & 0.887 & [0.875, 0.891] \\
    Power law on the gap                    & 3 & 0.033 & 0.012 & 0.005 & 0.006 & 0.889 & [0.881, 0.891] \\
    Loss-to-loss chain, logistic link       & 4 & 0.035 & 0.014 & 0.013 & 0.006 & 0.897 & [0.883, 0.901] \\
    Logistic in $\log C$                    & 4 & 0.038 & 0.017 & 0.014 & 0.006 & 0.883$^{\dagger}$ & [0.875, 0.886] \\
    BNSL, one break                         & 6 & 0.055 & 0.016 & 0.010 & 0.006 & 0.874$^{\dagger}$ & [0.871, 0.877] \\
    Two-stage via real NLL, exponential     & 3 & 0.071 & 0.023 & 0.006 & 0.018 & 1.000$^{\dagger}$ & [1.000, 1.000] \\
    BNSL, two breaks                        & 9 & ${>}0.3$ & 0.027 & 0.004 & 0.006 & 0.906 & [0.899, 0.914] \\
    \midrule
    M4, ceiling fixed at TabPFN's 0.887     & 3 & 0.015 & 0.006 & 0.002 & 0.005 & -- & -- \\
    \bottomrule
  \end{tabular}
\end{table}

\paragraph{Forms that interpolate equally well differ in how far they extrapolate.}
Leaving out one run at a time hardly separates the forms: all forms in Table~\ref{tab:forms} except the exponential
link reach an MAE of 0.005 to 0.007. Extrapolating over three decades or more does separate them. The best form with a
free ceiling is a power law in logit space with three parameters (0.024). The two-stage form through the real-data NLL
does about as well (0.022). The power law on the gap, M4 and the logistic in $\log C$ reach 0.03 to 0.04, and BNSL
with one break 0.055. Fitted on the points up to $3\times10^{13}$ or $10^{14}$ FLOPs, all of these forms err in the
same direction: they predict less improvement than the frontier later shows. On the frontier built from the five best
runs per level, the logit power law is again the best free form (0.039), next to the power law on the gap (0.040),
while the two-stage form falls behind (0.073; Appendix~\ref{app:functional-forms}). Fixing the ceiling at TabPFN's
0.887 removes about a third of the long-horizon error of the best free form on both frontiers. This shows how much an
external estimate of the ceiling is worth. It is not a test of the form, because the value was chosen after seeing the
data.

\paragraph{More breaks fit better and extrapolate worse.}
BNSL with two breaks has the best in-sample fit ($R^2=0.998$, against 0.991 for the power law on the gap) and the
smallest short-horizon error. Fitted on the points up to $10^{13}$ FLOPs, however, it diverges. Its four reduced variants and
the step forms do no better, with long-horizon errors from 0.05 to well above 0.3. A second break is supported only once both transitions lie inside the fitting range, which
our frontier does not provide early on.

\paragraph{Predicting through a loss does not make extrapolation more reliable.}
On the frontier over all runs, the two-stage forms with a logistic link reach long-horizon errors of 0.022 to 0.035,
comparable to the direct forms. On the top-5 frontier they reach 0.073 to 0.127, worse than the power laws. The
loss-to-loss chain, which predicts the real NLL from the synthetic validation loss, adds nothing to fitting the real NLL
directly. In both cases the fitted floor of the real NLL goes to zero, so the ceiling comes from the link alone. The
exponential link fits poorly ($R^2=0.75$) and sends the ceiling to 1.

\paragraph{The data do not determine the ceiling.}
On all 22 points, the five forms with the lowest long-horizon error put the ceiling between 0.887 and 0.909, around
TabPFN's 0.887 on the same tasks. For most forms the bootstrap interval is only 0.01 to 0.02 wide, but the intervals of
different forms do not all overlap: BNSL with one break gives $[0.871, 0.877]$, the logit power law $[0.885, 0.899]$.
More importantly, the estimate does not converge as the fitting range grows. Refitted on the points up to each compute
level, the ceilings of four representative forms stay 0.004 to 0.045 above the best ROC AUC observed so far from
$10^{13.75}$ FLOPs on, and rise with it, from 0.83--0.86 at $10^{14}$ FLOPs to 0.89--0.90 at $10^{16.5}$ FLOPs
(Table~\ref{tab:ceiling-drift}). The bootstrap
measures how much the ceiling varies with the runs, for a fixed fitting range. It does not measure how far the ceiling
will move once larger budgets are added.

\paragraph{Implications.}
For bounded downstream metrics of tabular foundation models, we recommend a three-parameter power law with an
irreducible term in logit space. It extrapolates best among the forms with a free ceiling on both frontiers, while
sigmoids, multi-break forms and two-stage forms do not. Forms should be ranked by long-horizon extrapolation on a
frontier built from all runs, with ceiling bounds and starting values taken from the training points only. The fitted
ceiling should be read as a lower bound that tracks the data, not as an estimate of irreducible performance. Matching
TabPFN's ROC AUC to three decimals is consistent with our data but not supported by it. These conclusions rest on one
set of 6{,}075 runs, trained with a constant learning rate, on 22 frontier points and on 21 datasets.
